\section{Estado del Arte y Mercado Potencial}

En los últimos años, el crecimiento sostenido de las grandes ciudades y el aumento del parque automotor impulsaron el desarrollo de múltiples soluciones tecnológicas orientadas a mejorar la movilidad urbana y asistir a los conductores durante sus desplazamientos. Dentro de este contexto, surgieron aplicaciones de navegación, sistemas de monitoreo de tránsito en tiempo real y plataformas de estacionamiento medido administradas por organismos municipales. Estas herramientas permiten optimizar recorridos, visualizar congestiones vehiculares y gestionar pagos asociados al uso del espacio público.

Sin embargo, gran parte de las soluciones actualmente disponibles se encuentran enfocadas principalmente en la circulación sobre la vía pública y en la administración de estacionamiento urbano regulado, dejando parcialmente desatendida la problemática relacionada con la disponibilidad, reserva y utilización eficiente de espacios privados de estacionamiento. En consecuencia, continúa existiendo una fragmentación significativa entre la oferta de cocheras privadas y la demanda de conductores que necesitan acceder rápidamente a un lugar seguro y cercano para estacionar.

A su vez, la creciente digitalización de servicios y la consolidación de modelos basados en economías colaborativas han demostrado la viabilidad de plataformas que permiten conectar usuarios con recursos ociosos mediante aplicaciones móviles y sistemas centralizados. Este paradigma abrió nuevas oportunidades para optimizar infraestructuras ya existentes sin necesidad de realizar grandes inversiones físicas adicionales.

En este escenario, ParkLink busca posicionarse como una plataforma capaz de integrar y centralizar distintos tipos de espacios de estacionamiento, permitiendo a conductores acceder a información unificada, realizar reservas anticipadas y gestionar pagos de manera digital y segura. El proyecto apunta especialmente a cubrir la falta de soluciones flexibles orientadas a reservas de corta duración y disponibilidad dinámica.

A nivel de mercado y competencia, pueden identificarse distintos modelos existentes que presentan limitaciones parciales frente a la propuesta planteada:

\begin{itemize}
\item \textbf{Estacionamiento libre y playas públicas:} representan la alternativa más tradicional utilizada por conductores urbanos. Sin embargo, su disponibilidad depende completamente de la ocupación inmediata, generando incertidumbre, tiempos prolongados de búsqueda y una fuerte variabilidad de precios según la zona y el horario.


\item \textbf{Sistemas de estacionamiento medido:} se encuentran limitados exclusivamente al uso de espacios públicos sobre la vía urbana. Generalmente no permiten reservas anticipadas y operan bajo restricciones horarias y regulatorias definidas por cada municipio.

\item \textbf{Plataformas de alquiler de cocheras tradicionales:} suelen estar orientadas a contratos mensuales o de larga duración, priorizando usuarios permanentes por sobre necesidades temporales o transaccionales. Esto reduce considerablemente la flexibilidad para conductores que requieren estacionamiento únicamente por horas o períodos breves.

\item \textbf{Aplicaciones de navegación y tráfico:} herramientas como sistemas GPS o plataformas de mapas colaborativos brindan información sobre tránsito y rutas, pero no ofrecen integración completa con reservas de estacionamiento privado ni mecanismos centralizados de contratación.


\end{itemize}

A partir de este análisis, se observa una oportunidad concreta para el desarrollo de una plataforma que combine geolocalización, reserva anticipada, gestión digital de pagos y utilización colaborativa de espacios subutilizados. De esta manera, ParkLink busca diferenciarse mediante un enfoque centrado en la flexibilidad, accesibilidad y optimización de recursos urbanos existentes, apuntando inicialmente a contextos urbanos de alta densidad como la Ciudad Autónoma de Buenos Aires.
